\begin{theorem}[escort 温度族的两态 Gibbs 化与末位统计的渐近充分性]\label{thm:killo-fold-bin-escort-onebit-fisher}
对每个 \(m\ge 1\)，记
\[
A_0:=\{w\in X_m:\ w_m=0\},
\qquad
A_1:=\{w\in X_m:\ w_m=1\},
\]
并令 \(u_m\) 为 \(X_m\) 上的均匀分布。
对任意 \(q\ge 0\)，定义 escort 分布
\[
\pi_{m,q}(w):=\frac{\bigl(d_m^{\mathrm{bin}}(w)\bigr)^q}{S_q^{\mathrm{bin}}(m)},
\qquad
S_q^{\mathrm{bin}}(m):=\sum_{x\in X_m}\bigl(d_m^{\mathrm{bin}}(x)\bigr)^q,
\]
以及两态 Gibbs 近似
\[
\nu_{m,q}(w):=\frac{u_m(w)\varphi^{-q w_m}}{Z_{m,q}},
\qquad
Z_{m,q}:=\frac{F_{m+1}+\varphi^{-q}F_m}{F_{m+2}}.
\]
则对任意有界温度区间 \(q\in[0,Q]\)，存在常数 \(C_Q>0\) 使对所有充分大的 \(m\) 成立
\[
\sup_{q\in[0,Q]}\sup_{w\in X_m}
\left|
\frac{\pi_{m,q}(w)}{\nu_{m,q}(w)}-1
\right|
\le C_Q\Bigl(\frac{\varphi}{2}\Bigr)^m.
\]
从而
\[
\sup_{q\in[0,Q]}
\|\pi_{m,q}-\nu_{m,q}\|_{\mathrm{TV}}
=
O_Q\!\Bigl(\Bigl(\frac{\varphi}{2}\Bigr)^m\Bigr).
\]
特别地，末位映射 \(w\mapsto w_m\in\{0,1\}\) 构成 escort 温度族的渐近充分统计量：对 \(i\in\{0,1\}\)，有
\[
\sup_{q\in[0,Q]}
\bigl\|
\pi_{m,q}(\,\cdot\,|A_i)-u_m(\,\cdot\,|A_i)
\bigr\|_{\mathrm{TV}}
=
O_Q\!\Bigl(\Bigl(\frac{\varphi}{2}\Bigr)^m\Bigr).
\]
\end{theorem}
\begin{proof}
由定理 \ref{thm:fold-bin-two-state-asymptotic} 与命题 \ref{prop:fold-bin-power-sum-asymptotic}，对每个有界区间 \(q\in[0,Q]\) 都有一致展开
\[
\bigl(d_m^{\mathrm{bin}}(w)\bigr)^q
=
\Bigl(\frac{2}{\varphi}\Bigr)^{qm}\varphi^{-q w_m}
\Bigl(1+O_Q\!\Bigl(\Bigl(\frac{\varphi}{2}\Bigr)^m\Bigr)\Bigr),
\qquad
(w\in X_m),
\]
以及
\[
S_q^{\mathrm{bin}}(m)
=
\Bigl(\frac{2}{\varphi}\Bigr)^{qm}
\bigl(F_{m+1}+\varphi^{-q}F_m\bigr)
\Bigl(1+O_Q\!\Bigl(\Bigl(\frac{\varphi}{2}\Bigr)^m\Bigr)\Bigr),
\]
其中误差对 \(q\in[0,Q]\) 与 \(w\in X_m\) 一致。
将两式相除，并用 \(u_m(w)=1/F_{m+2}\)，得到
\[
\pi_{m,q}(w)
=
\frac{\varphi^{-q w_m}}{F_{m+1}+\varphi^{-q}F_m}
\Bigl(1+O_Q\!\Bigl(\Bigl(\frac{\varphi}{2}\Bigr)^m\Bigr)\Bigr)
=
\nu_{m,q}(w)
\Bigl(1+O_Q\!\Bigl(\Bigl(\frac{\varphi}{2}\Bigr)^m\Bigr)\Bigr),
\]
即得逐点乘法逼近。

由 \(\sum_{w\in X_m}\nu_{m,q}(w)=1\)，直接得到
\[
\|\pi_{m,q}-\nu_{m,q}\|_{\mathrm{TV}}
\le
\frac12\sum_{w\in X_m}\nu_{m,q}(w)\,
O_Q\!\Bigl(\Bigl(\frac{\varphi}{2}\Bigr)^m\Bigr)
=
O_Q\!\Bigl(\Bigl(\frac{\varphi}{2}\Bigr)^m\Bigr).
\]

再看条件分布。
由定义，\(\nu_{m,q}\) 在每个块 \(A_i\) 上严格均匀；并且
\[
\nu_{m,q}(A_1)=\frac{\varphi^{-q}F_m}{F_{m+1}+\varphi^{-q}F_m},
\qquad
\nu_{m,q}(A_0)=\frac{F_{m+1}}{F_{m+1}+\varphi^{-q}F_m}.
\]
由于 \(q\in[0,Q]\) 且 \(F_{m+1}/F_m\to\varphi\)，存在常数 \(c_Q>0\) 使对所有充分大的 \(m\) 有
\[
c_Q\le \nu_{m,q}(A_i)\le 1-c_Q
\qquad
(i=0,1,\ q\in[0,Q]).
\]
另一方面，由逐点乘法逼近，对 \(w\in A_i\) 有
\[
\pi_{m,q}(w)=\nu_{m,q}(w)\Bigl(1+O_Q\!\Bigl(\Bigl(\frac{\varphi}{2}\Bigr)^m\Bigr)\Bigr),
\qquad
\pi_{m,q}(A_i)=\nu_{m,q}(A_i)\Bigl(1+O_Q\!\Bigl(\Bigl(\frac{\varphi}{2}\Bigr)^m\Bigr)\Bigr).
\]
因此
\[
\pi_{m,q}(w|A_i)
=
\frac{\nu_{m,q}(w)}{\nu_{m,q}(A_i)}
\Bigl(1+O_Q\!\Bigl(\Bigl(\frac{\varphi}{2}\Bigr)^m\Bigr)\Bigr)
=
\nu_{m,q}(w|A_i)
\Bigl(1+O_Q\!\Bigl(\Bigl(\frac{\varphi}{2}\Bigr)^m\Bigr)\Bigr).
\]
而 \(\nu_{m,q}(\cdot|A_i)=u_m(\cdot|A_i)\) 在 \(A_i\) 上均匀，故
\[
\bigl\|
\pi_{m,q}(\,\cdot\,|A_i)-u_m(\,\cdot\,|A_i)
\bigr\|_{\mathrm{TV}}
=
O_Q\!\Bigl(\Bigl(\frac{\varphi}{2}\Bigr)^m\Bigr),
\]
即证。
\end{proof}

\leanverified{paper\_killo\_fold\_bin\_escort\_onebit\_fisher}

\leanverified{paper\_killo\_fold\_bin\_escort\_renyi\_logistic\_geometry}
\begin{theorem}[escort 温度族的全阶 Rényi 信息几何由 logistic 曲线控制]\label{thm:killo-fold-bin-escort-renyi-logistic-geometry}
定义
\[
\theta(q):=\frac{1}{1+\varphi^{q+1}},
\qquad q\ge 0.
\]
则对任意固定 \(p,q\ge 0\) 与任意固定 \(\alpha\in(0,\infty)\setminus\{1\}\)，有
\[
\lim_{m\to\infty}D_\alpha(\pi_{m,q}\Vert \pi_{m,p})
=
\frac{1}{\alpha-1}
\log\!\Bigl(
\theta(q)^\alpha\theta(p)^{1-\alpha}
+
\bigl(1-\theta(q)\bigr)^\alpha\bigl(1-\theta(p)\bigr)^{1-\alpha}
\Bigr).
\]
当 \(\alpha\to 1\) 时，上式连续退化为
\[
\lim_{m\to\infty}D_{\mathrm{KL}}(\pi_{m,q}\Vert \pi_{m,p})
=
\theta(q)\log\frac{\theta(q)}{\theta(p)}
+
\bigl(1-\theta(q)\bigr)\log\frac{1-\theta(q)}{1-\theta(p)}.
\]
此外，函数 \(q\mapsto \theta(q)\) 满足 logistic 微分方程
\[
\theta'(q)=-(\log\varphi)\,\theta(q)\bigl(1-\theta(q)\bigr),
\]
从而其 Fisher 信息为
\[
I(q):=\frac{\theta'(q)^2}{\theta(q)(1-\theta(q))}
=(\log\varphi)^2\theta(q)\bigl(1-\theta(q)\bigr).
\]
并且 \(q\mapsto I(q)\) 在 \([0,\infty)\) 上严格递减。
\end{theorem}
\begin{proof}
记
\[
\theta_m(q):=\frac{\varphi^{-q}F_m}{F_{m+1}+\varphi^{-q}F_m}.
\]
由定理 \ref{thm:killo-fold-bin-escort-onebit-fisher}，对固定 \(p,q\) 有
\[
\pi_{m,q}(w)=\nu_{m,q}(w)\bigl(1+O((\varphi/2)^m)\bigr),
\qquad
\pi_{m,p}(w)=\nu_{m,p}(w)\bigl(1+O((\varphi/2)^m)\bigr),
\]
其中误差在 \(w\in X_m\) 上一致。
因此对固定 \(\alpha\in(0,\infty)\setminus\{1\}\)，有
\[
\pi_{m,q}(w)^\alpha\pi_{m,p}(w)^{1-\alpha}
=
\nu_{m,q}(w)^\alpha\nu_{m,p}(w)^{1-\alpha}
\Bigl(1+O\!\Bigl(\Bigl(\frac{\varphi}{2}\Bigr)^m\Bigr)\Bigr),
\]
并且该误差仍对 \(w\) 一致。
于是
\[
\sum_{w\in X_m}\pi_{m,q}(w)^\alpha\pi_{m,p}(w)^{1-\alpha}
=
\Bigl(\sum_{w\in X_m}\nu_{m,q}(w)^\alpha\nu_{m,p}(w)^{1-\alpha}\Bigr)
\Bigl(1+O\!\Bigl(\Bigl(\frac{\varphi}{2}\Bigr)^m\Bigr)\Bigr).
\]

另一方面，\(\nu_{m,q}\) 在两块上的表达为
\[
\nu_{m,q}(w)=
\begin{cases}
\dfrac{1-\theta_m(q)}{|A_0|}, & w\in A_0,\\[2ex]
\dfrac{\theta_m(q)}{|A_1|}, & w\in A_1.
\end{cases}
\]
于是
\[
\sum_{w\in X_m}\nu_{m,q}(w)^\alpha\nu_{m,p}(w)^{1-\alpha}
=
\bigl(1-\theta_m(q)\bigr)^\alpha\bigl(1-\theta_m(p)\bigr)^{1-\alpha}
+
\theta_m(q)^\alpha\theta_m(p)^{1-\alpha},
\]
其中块大小完全相消。
故
\[
D_\alpha(\nu_{m,q}\Vert \nu_{m,p})
=
\frac{1}{\alpha-1}
\log\!\Bigl(
\theta_m(q)^\alpha\theta_m(p)^{1-\alpha}
+
\bigl(1-\theta_m(q)\bigr)^\alpha\bigl(1-\theta_m(p)\bigr)^{1-\alpha}
\Bigr).
\]
由上一式可知
\[
D_\alpha(\pi_{m,q}\Vert \pi_{m,p})
=
D_\alpha(\nu_{m,q}\Vert \nu_{m,p})+o(1),
\]
再由 \(\theta_m(\cdot)\to \theta(\cdot)\) 即得所示 Rényi 极限。

当 \(\alpha\to 1\) 时，Rényi 散度向 KL 散度连续延拓，得到
\[
\theta(q)\log\frac{\theta(q)}{\theta(p)}
+
\bigl(1-\theta(q)\bigr)\log\frac{1-\theta(q)}{1-\theta(p)},
\]
这与命题 \ref{prop:fold-bin-escort-escort-kl} 一致。

最后，由
\[
\theta(q)=\frac{1}{1+\varphi^{q+1}}
\]
直接求导得
\[
\theta'(q)
=
-\frac{\varphi^{q+1}\log\varphi}{(1+\varphi^{q+1})^2}
=
-(\log\varphi)\,\theta(q)\bigl(1-\theta(q)\bigr).
\]
因此
\[
I(q)=\frac{\theta'(q)^2}{\theta(q)(1-\theta(q))}
=(\log\varphi)^2\theta(q)\bigl(1-\theta(q)\bigr).
\]
又因
\[
0<\theta(q)\le \theta(0)=\frac{1}{1+\varphi}=\varphi^{-2}<\frac12,
\]
且 \(q\mapsto \theta(q)\) 严格递减，而 \(x\mapsto x(1-x)\) 在 \((0,1/2]\) 上严格递增，故
\[
q\mapsto \theta(q)\bigl(1-\theta(q)\bigr)
\]
在 \([0,\infty)\) 上严格递减，从而 \(I(q)\) 亦严格递减。
\end{proof}

\endinput
